\section{Ablation}
\label{sec:ablation}

The proposed solver has five components: the cap on the scaled rates of \cref{eq:scaled}, the refusal rule of \cref{eq:refuse-horizon,eq:refuse-base,eq:refuse-unreachable}, the continuous window function of \cref{eq:window}, the upward rounding with the neighbour check of \cref{eq:rounding}, and the effective capacity that folds the inflows and the flows outside the segments into $C_t$. \Cref{tab:ablation} removes one at a time on the \kInst{} instances of the corpus and scores the variant against the oracle with the metrics of \cref{eq:metrics}.

\begin{table*}[htbp]
\centering
\small
\setlength{\tabcolsep}{3pt}
\caption{Ablation on the corpus of 20,000 instances: each row removes one component of the proposed solver; the columns are the share of refused instances answered with a number, the share of solved instances left without a number, the share of solved instances whose multiplier is a threshold, the share of all instances answered with a multiplier above the cap, and the median relative error.}
\label{tab:ablation}
\begin{tabular}{p{4.4cm}p{1.9cm}p{1.9cm}p{1.9cm}p{1.9cm}p{1.9cm}}
\toprule
Variant & answer if refused (\%) & no answer if solved (\%) & threshold (\%) & beyond cap (\%) & median rel. error \\
\midrule
proposed solver (all components) & 0.0 & 0.0 & 100.0 & 0.00 & 8.9e-08 \\
without the cap on the scaled rates & 4.6 & 0.0 & 65.7 & 2.04 & 2.3e-07 \\
without the refusal rule & 100.0 & 0.0 & 99.0 & 3.48 & 9.2e-08 \\
without the continuous window function (Brent on SP) & 13.5 & 0.0 & 15.7 & 0.00 & 6.3e-02 \\
without upward rounding and the neighbour check & 0.0 & 0.0 & 51.1 & 0.00 & 4.6e-08 \\
without inflows and flows outside the segments & 1.4 & 4.6 & 34.4 & 0.00 & 1.5e-02 \\
\bottomrule
\end{tabular}
\end{table*}


Removing the cap leaves the solver able to answer where no admissible multiplier exists. The variant searches up to a multiplier of 1,000 with rates $m\,q_{j,d}$ that may exceed 1, and a rate above 1 turns the stock factor $1-q$ negative, so the stock changes sign and the cumulative outflow is no longer a share of the opening balance. It answers \kUFalse{}\,\% of the refused instances, and it returns a multiplier above the cap on \kUBeyond{}\,\% of all instances. Its multipliers are not thresholds of the model: they satisfy the threshold property on \kUThr{}\,\% of the solved instances, and the experiment did not isolate which of the artefacts of rates above 1 causes the failures. The cap is not a numerical convenience. It is what makes the scaled scenario a run-off scenario.

Removing the refusal rule leaves a bisection on a fixed bracket that always answers. It answers all \kRefused{} refused instances with an endpoint of its bracket, and on \kNBeyond{}\,\% of all instances the endpoint lies above the cap of the scenario, so the number it reports cannot be reached by any admissible run. Where it solves, it is as accurate as the proposed solver (median relative error \kNErr{}) and satisfies the threshold property on \kNThr{}\,\% of the solved instances; the cost is \kNEvals{} evaluations against \kPEvals{}, because a bisection does not use the shape of the function. The component is therefore not about accuracy. It decides whether a number is a result.

Removing the continuous window function means applying Brent's method to the integer survival period. The variant cannot bracket a target that the base already reaches or that no multiplier reaches, since $\mathrm{SP}(m)-T$ then has no sign change, and it answers \kBFalse{}\,\% of the refused instances where a sign change happens to exist. On the solved instances it returns a threshold on \kBThr{}\,\% only, because it stops inside the interval on which the survival period equals the target; its median error is \kBErrPc{}\,\%, and its 95th percentile \kBErrQ{}. The window function removes exactly that flat region.

Removing the upward rounding and the neighbour check, so that the root is rounded to the nearest decimal, keeps the refusals and the accuracy, with a median error of \kRErr{}, and fails the threshold property on about half of the solved instances: it satisfies it on \kRThr{}\,\%. The reason is the one in \cref{prop:threshold}: the set of multipliers that reach the target is open at the root, so a rounded-down value does not reach it. The check is cheap, and it gives the reader a number that reproduces.

Removing the effective capacity hands the solver the capacity schedule alone. It answers \kWFalse{}\,\% of the refused instances and fails to answer \kWMiss{}\,\% of the solved ones, since the position it solves is not the position of the ladder. On the instances that have inflows or flows outside the segments (\kRawFlowsN{} solved ones) the median error of its multiplier is \kRawFlows{}\,\%, and on the committed bank it moves the Adverse multiplier from \mAdv{} to \rawAdv{} and the Extreme one from \mExt{} to \rawExt{}. A reverse stress that is solved on a position other than the ladder's is a reverse stress of another bank.

Each component protects against a different failure: the cap against infeasible scenarios, the refusal rule against numbers for undefined targets, the window function against flat regions of the survival period, the rounding against a multiplier that does not reach the target, and the effective capacity against a position that differs from the ladder's. No single component carries the result; the proposed solver is the only variant at zero false answers, zero misses, no multiplier above the cap and the threshold property on every solved instance.
